% 文件开头添加条件判断
\newif\ifstandalone
\standalonetrue  % 默认独立编译模式

% 检测是否在主文档中
\ifdefined\maindoc
  \standalonefalse
\fi

\ifstandalone
  \documentclass[12pt,a4paper]{ctexart}

  \usepackage[utf8]{inputenc}
  \usepackage{amsmath,amssymb,amsthm}
  \usepackage{graphicx}
  \usepackage{hyperref}
  \usepackage{geometry}
  \usepackage{array}
  \usepackage{booktabs}
  \geometry{a4paper, margin=2.5cm}

  \title{第九章\quad 力控制}
  \author{Robot Manipulator Control - 中文翻译}
  \date{\today}

  \newtheorem{definition}{定义}[section]
\fi

\ifstandalone
  \begin{document}
  \maketitle
\fi

本章研究位置/力控制器的基本原理。涵盖的主题包括刚度控制、混合位置/力控制、混合阻抗控制和降维状态位置/力控制。重点放在控制器的开发、稳定性和实现问题上。

%===========================================
\section{引言}\label{sec:9.1}
%===========================================

对于机器人机械臂执行的任务，例如移动载荷或喷涂物体，位置控制器能够提供足够的性能，因为这些类型的任务仅要求机器人跟随期望轨迹。然而，在磨削或装配任务中，机器人机械臂与环境接触；因此，在机器人机械臂和环境之间会产生相互作用力。因此，这些相互作用力以及末端执行器的位置都必须被控制。

为了说明同时使用力和位置控制的必要性，考虑控制机械臂在黑板上书写句子的问题。为了形成句子中的字母，我们当然必须控制末端执行器的位置，或者等价地，控制粉笔的位置。正如任何在黑板上书写过的人都知道的，按压黑板的力也必须被控制。也就是说，按压太轻会导致字母不易辨认，而按压太重会导致粉笔折断。这个例子清楚地说明了许多机器人应用将需要指定期望位置轨迹和期望力轨迹。在本章中，我们介绍一些通用的控制策略，不仅控制机器人末端执行器的位置，还控制末端执行器施加在环境上的力。需要注意的是，在本章中，我们假设控制器发出的期望速度和力轨迹与环境模型一致[Lipkin and Duffy 1988]。如果不是这种情况，可以修改期望速度和力轨迹以使其与环境模型一致。有兴趣的读者可以参考[Lipkin and Duffy 1988]了解关于这种修改或期望轨迹的``运动静力学滤波''的信息。

%===========================================
\section{刚度控制}\label{sec:9.2}
%===========================================

由于工业过程中涉及的第一批机器人机械臂需要执行位置任务（例如，喷涂），机器人机械臂被制造得非常刚性。这种刚性设计允许机器人控制设计者通过使用简单的控制律获得合理的位置精度。正如人们所预期的，力控制应用（例如，磨削或打磨）用这种``刚性''机器人极其难以完成。因此，如果机器人机械臂的``刚度''可以被控制，力控制应用将更容易实现。在本节中，针对简单的单自由度示例阐述了刚度控制的概念。然后修改机器人机械臂方程以考虑施加在环境上的力。使用这个新模型，刚度控制概念[Salisbury and Craig 1980]被推广到$n$连杆机器人机械臂。

%-------------------------------------------
\subsection{单自由度机械臂的刚度控制}
%-------------------------------------------

为了阐述刚度控制的概念，考虑图9.2.1所示系统的力控制问题。这里质量为$m$的机械臂被假设与环境接触，环境位于静态位置$x_e$。控制问题是确定输入力（即$\tau$），使机械臂移动到期望的恒定位置（即$x_d$）。在这个系统中，我们还假设如果机械臂的位置（即$x$）大于$x_e$，施加在环境上的力（即$f$）由下式给出：
\begin{equation}\label{eq:9.2.1}
f = k_e(x - x_e),
\end{equation}
其中$k_e$是一个正常数，用于表示环境刚度。也就是说，我们假设环境刚度可以被建模为弹簧常数为$k_e$的线性弹簧。基于这个假设，我们可以将单自由度系统可视化为图9.2.2所示的质量-弹簧图。

假设重力和摩擦可以忽略不计，图9.2.2所示系统的运动方程由下式给出：
\begin{equation}\label{eq:9.2.2}
m\ddot{x} = \tau - f
\end{equation}
式(\ref{eq:9.2.2})的系统框图由图9.2.3给出。注意，在图9.2.3中，我们使用变量$s$表示拉普拉斯变换变量。

式(\ref{eq:9.2.2})的动力学形式启发了简单的PD控制律：
\begin{equation}\label{eq:9.2.3}
\tau = -k_v\dot{x} + k_p(x_d - x)
\end{equation}
其中$k_v$和$k_p$是正标量控制增益。将式(\ref{eq:9.2.3})代入式(\ref{eq:9.2.2})后，我们得到闭环系统：
\begin{equation}\label{eq:9.2.4}
m\ddot{x} + k_v\dot{x} + k_p x + k_e(x - x_e) = k_p x_d
\end{equation}
这可以用图9.2.4所示的框图表示。从图9.2.4我们知道闭环系统是稳定的，因为我们已经定义常数$m$、$k_v$、$k_p$和$k_e$为正数。也就是说，我们可以证明传递函数
\begin{equation}\label{eq:9.2.transfer}
H(s) = \frac{1}{ms^2 + k_v s + (k_p + k_e)}
\end{equation}
的极点在开左半$s$平面。

为了研究式(\ref{eq:9.2.3})给出的PD控制如何控制施加在环境上的力，我们检查系统在稳态条件下的情况。因为$x_d$和$x_e$是常数，可以从式(\ref{eq:9.2.4})找到$x$的拉普拉斯变换为
\begin{equation}\label{eq:9.2.5}
x(s) = \frac{k_p x_d + k_e x_e}{s(ms^2 + k_v s + (k_p + k_e))}
\end{equation}
因此，稳态机械臂位置（即$\bar{x}$）可以很容易地证明为
\begin{equation}\label{eq:9.2.6}
\bar{x} = \lim_{s \to 0} s\,x(s) = \frac{k_p x_d + k_e x_e}{k_p + k_e}
\end{equation}

现在可以使用稳态机械臂位置来计算施加在环境上的稳态力（即$\bar{f}$）。具体来说，将式(\ref{eq:9.2.6})代入式(\ref{eq:9.2.1})，稳态力由下式给出：
\begin{equation}\label{eq:9.2.7}
\bar{f} = \frac{k_p k_e(x_d - x_e)}{k_p + k_e}
\end{equation}

正如人们所预期的，环境弹簧常数通常被认为很大，因为在大多数机器人力控制应用中，机器人推动的是几乎刚性的表面。因此，如果我们假设$k_e \gg k_p$，我们可以将式(\ref{eq:9.2.7})中的稳态力近似为
\begin{equation}\label{eq:9.2.8}
\bar{f} \approx k_p(x_d - x_e)
\end{equation}

从上面的讨论中，我们可以看到式(\ref{eq:9.2.3})给出的位置控制策略确实会对环境施加一个力。具体来说，这个力是通过命令一个稍微进入接触面内部的期望轨迹而产生的。在试图消除位置误差时，位置控制器在表面上施加了一个稳态力。从式(\ref{eq:9.2.8})给出的近似稳态力来看，位置增益（即$k_p$）可以被认为是代表机械臂的期望``刚度''。也就是说，机械臂可以被可视化为一个弹簧，弹簧常数为$k_p$，对环境施加力。因此，``刚度控制''这个术语经常与式(\ref{eq:9.2.3})给出的PD控制相关联，因为机械臂的刚度可以通过调整$k_p$来设置。

%-------------------------------------------
\subsection{雅可比矩阵和环境力}
%-------------------------------------------

在将刚度控制器推广到$n$连杆机器人机械臂之前，我们必须定义一些关于机器人对环境施加的力的符号。正如[Spong and Vidyasagar 1989]中所解释的，力通常通过雅可比矩阵转换到关节空间。在本章中，我们根据为特定机器人应用定义的任务空间坐标系来定义雅可比矩阵。也就是说，对于某个应用，我们可能希望末端执行器沿一组特定方向施加力，而沿其他方向移动。这个概念由图9.2.5说明，图中描述了一个沿倾斜表面移动的机械臂。任务空间坐标系由方向$u$和$v$给出，因为我们希望沿$v$方向沿表面移动末端执行器，同时沿$u$方向对表面施加法向力。通常从这种推理中为大多数机器人应用定义适当的任务空间坐标系。

按照这个逻辑，让$x$为由下式定义的$n \times 1$任务空间向量
\begin{equation}\label{eq:9.2.9}
x = h(q)
\end{equation}
其中$h(q)$从机械臂运动学和关节空间与任务空间之间的适当关系中找到。$x$的导数定义为
\begin{equation}\label{eq:9.2.10}
\dot{x} = J(q)\dot{q}
\end{equation}
其中$n \times n$任务空间雅可比矩阵$J(q)$[Spong and Vidyasagar 1989]定义为
\begin{equation}\label{eq:9.2.11}
J(q) = \begin{bmatrix} I & 0 \\ 0 & T \end{bmatrix} \frac{\partial h(q)}{\partial q}
\end{equation}
其中单位矩阵$I$、零矩阵$0$和变换矩阵$T$的维度取决于所选的任务空间坐标系。变换矩阵$T$通常在将关节速度转换为与末端执行器方向相关的横滚、俯仰和偏航角的导数时使用。为简洁起见，我们假设本章讨论的机器人机械臂是非冗余的，并且始终处于非奇异构型；因此，雅可比矩阵是一个非奇异方阵。

使用任务空间坐标概念，我们现在检查如何修改机器人方程以用于力控制目的。如果机械臂的末端执行器与环境接触，末端执行器和环境之间将存在力相互作用。如果相互作用力在关节空间中测量，机械臂动力学方程可以写成
\begin{equation}\label{eq:9.2.12}
\tau = M(q)\ddot{q} + V_m(q, \dot{q})\dot{q} + G(q) + F(\dot{q}) + \tau_e
\end{equation}
其中$\tau_e$是关节空间坐标中的$n \times 1$向量，表示施加在环境上的力。式(\ref{eq:9.2.12})给出的动力学方程是有意义的，因为如果机械臂不运动（即$\dot{q} = \ddot{q} = 0$），则式(\ref{eq:9.2.12})简化为
\begin{equation}\label{eq:9.2.13}
\tau = G(q) + \tau_e
\end{equation}
也就是说，如果机器人机械臂处于静态操作，执行器力等于施加在环境上的力加上承受重力所需的力。注意，在式(\ref{eq:9.2.13})中，我们假设可以忽略静摩擦。

关节空间表示施加在环境上的力不是机器人文献中的标准符号；相反，机器人机械臂方程通常由下式给出
\begin{equation}\label{eq:9.2.14}
\tau = M(q)\ddot{q} + V_m(q, \dot{q})\dot{q} + G(q) + F(\dot{q}) + J^T(q)f
\end{equation}
其中$f$是任务空间中的接触力和力矩的$n \times 1$向量。

为了理解式(\ref{eq:9.2.14})的来源，将式(\ref{eq:9.2.12})和式(\ref{eq:9.2.14})的右侧相等，得到
\begin{equation}\label{eq:9.2.15}
\tau_e = J^T(q)f
\end{equation}
这个方程可以通过使用能量守恒论证很容易地证明为真。具体来说，通过能量守恒，我们知道
\begin{equation}\label{eq:9.2.16}
\dot{q}^T \tau_e = \dot{x}^T f
\end{equation}
将式(\ref{eq:9.2.10})代入式(\ref{eq:9.2.16})得到
\[
\dot{q}^T \tau_e = (J(q)\dot{q})^T f = \dot{q}^T J^T(q)f
\]
由于上述关系必须对所有$q$都成立，我们可以看到式(\ref{eq:9.2.15})显然代表一个真命题。为了阐明开发雅可比矩阵和任务空间公式的过程，现在讨论两个示例。

\textbf{示例 9.2--1：倾斜表面的任务空间公式}

我们希望找到图9.2.5中所示笛卡尔机械臂系统（即两个关节都是棱柱形）的机械臂动力学，并将施加在表面上的力分解为法向力和切向力。首先，当机器人不受表面约束时，可以轻松确定动力学部分。在移除表面和相互作用力$f_1$和$f_2$后，机械臂动力学可以表示为
\begin{equation}
\tau = M\ddot{q} + G + F(\dot{q})
\end{equation}
其中
\[
M = \begin{bmatrix} m_1 & 0 \\ 0 & m_1 + m_2 \end{bmatrix}, \quad q = \begin{bmatrix} q_1 \\ q_2 \end{bmatrix}, \quad \tau = \begin{bmatrix} \tau_1 \\ \tau_2 \end{bmatrix}, \quad G = \begin{bmatrix} 0 \\ (m_1+m_2)g \end{bmatrix}
\]
$F(\dot{q})$是$2 \times 1$向量$\begin{bmatrix} F_1(\dot{q}_1) \\ F_2(\dot{q}_2) \end{bmatrix}$，它模拟了如第2章所讨论的摩擦。

为了考虑相互作用力，让$x$为由下式定义的$2 \times 1$任务空间向量
\begin{equation}
x = \begin{bmatrix} u \\ v \end{bmatrix}
\end{equation}
其中$u$和$v$定义了一个固定坐标系，使得$u$表示到表面的法向距离，$v$表示沿表面的切向距离。如式(\ref{eq:9.2.9})所示，任务空间坐标可以用关节空间坐标表示为
\begin{equation}
x = h(q)
\end{equation}
其中$h(q)$从问题的几何中找到为
\begin{equation}
h(q) = \frac{1}{\sqrt{2}} \begin{bmatrix} q_1 - q_2 \\ q_1 + q_2 \end{bmatrix}
\end{equation}

任务空间雅可比矩阵从式(\ref{eq:9.2.11})中找到，利用$T$是单位矩阵的事实，因为我们不需要关心任何末端执行器方向角。也就是说，$J(q)$由下式给出
\begin{equation}
J(q) = \frac{1}{\sqrt{2}} \begin{bmatrix} 1 & -1 \\ 1 & 1 \end{bmatrix}
\end{equation}

根据式(\ref{eq:9.2.14})，机器人机械臂方程由下式给出
\begin{equation}
M\ddot{q} + G + F(q) + J^T f = \tau
\end{equation}
其中
\[
G = \begin{bmatrix} 0 \\ (m_1+m_2)g \end{bmatrix}, \quad f = \begin{bmatrix} f_1 \\ f_2 \end{bmatrix}
\]
重要的是要认识到法向力（即$f_1$）和切向力（即$f_2$）是按(2)给出的任务空间坐标系的方向绘制的（见图9.2.5）。

\textbf{示例 9.2--2：椭圆表面的任务空间公式}

我们希望找到图9.2.6中所示笛卡尔机械臂系统的机械臂动力学，并将施加在表面上的力分解为法向力和切向力。动力学部分与示例9.2.1相同；然而，由于环境表面的变化，必须定义一个新的任务空间坐标系。具体来说，让$x$为由下式定义的$2 \times 1$任务空间向量
\begin{equation}
x = \begin{bmatrix} u \\ v \end{bmatrix}
\end{equation}
其中$u$和$v$定义了一个旋转坐标系，使得$u$表示到表面的法向距离，$v$表示沿表面的切向距离。如式(\ref{eq:9.2.9})所示，任务空间坐标可以用关节空间坐标表示为
\begin{equation}
x = h(q)
\end{equation}
其中$h(q)$求得为
\begin{equation}\tag{3}
h(q) = \begin{bmatrix} \bar{u} \cdot \bar{q}_1 & \bar{u} \cdot \bar{q}_2 \\ \bar{\nu} \cdot \bar{q}_1 & \bar{\nu} \cdot \bar{q}_2 \end{bmatrix} \begin{bmatrix} q_1 \\ q_2 \end{bmatrix}
\end{equation}

式(3)中的$\bar{u}$、$\bar{\nu}$分别为沿任务空间法向与切向的单位向量；$\bar{q}_1$、$\bar{q}_2$为沿关节坐标轴方向的固定单位基向量（勿与关节位移标量$q_1$、$q_2$混淆）。这些基向量取为
\begin{equation}\tag{4}
\bar{q}_1 = \begin{bmatrix} 1 \\ 0 \end{bmatrix}, \quad \bar{q}_2 = \begin{bmatrix} 0 \\ 1 \end{bmatrix}
\end{equation}

为了找到单位向量$\bar{u}$和$\bar{\nu}$，我们首先使用表面函数
\begin{equation}
\frac{1}{4}q_1^2 + q_2^2 = 1
\end{equation}
用一个变量（即$q_2$）参数化表面如下：
\begin{equation}\tag{5}
q_1 = 2\sqrt{1 - q_2^2}
\end{equation}

(6)关于$q_2$的偏导数给出沿参数曲线的切向方向；将其归一化后得到与表面相切的单位向量$\bar{\nu}$：
\begin{equation}\tag{7}
\bar{\nu} = \frac{1}{\Delta}
\begin{bmatrix} \dfrac{\partial q_1}{\partial q_2} \\ \dfrac{\partial q_2}{\partial q_2} \end{bmatrix}
= \frac{1}{\Delta}
\begin{bmatrix} -2q_2(1-q_2^2)^{-1/2} \\ 1 \end{bmatrix},
\end{equation}
\begin{equation*}
\Delta = \sqrt{1 + 4q_2^2(1-q_2^2)^{-1}} \, .
\end{equation*}

再次利用(5)，$\bar{\nu}$的表达式可简化为
\begin{equation}\tag{8}
\bar{\nu} = \frac{1}{\Delta}
\begin{bmatrix} -4q_2/q_1 \\ 1 \end{bmatrix},
\end{equation}
\begin{equation*}
\Delta = \sqrt{1 + 16q_2^2/q_1^2} \, .
\end{equation*}

因为向量$\bar{u}$和$\bar{\nu}$必须正交（即$\bar{u} \cdot \bar{\nu} = 0$），由(8)及几何关系可取（(8)中的$\Delta$与下文相同）
\begin{equation}\tag{9}
\bar{u} = \frac{1}{\Delta} \begin{bmatrix} 1 \\ 4q_2/q_1 \end{bmatrix}
\end{equation}

将(4)、(8)和(9)代入(3)得到
\begin{equation}\tag{10}
h(q) = \frac{1}{\Delta} \begin{bmatrix} q_1 + 4q_2^2/q_1 \\ -3q_2 \end{bmatrix}
\end{equation}

任务空间雅可比矩阵从式(\ref{eq:9.2.11})中找到，利用$T$是单位矩阵的事实。也就是说，$J(q)=\partial h/\partial q^{\mathsf T}$为
\begin{equation}\tag{11}
J(q) = \begin{bmatrix} J_{11} & J_{12} \\ J_{21} & J_{22} \end{bmatrix}
\end{equation}
其中
\begin{align*}
J_{11} &= \frac{1}{2q_1}\left(\frac{1}{4}q_1^2 + 7q_2^2\right)\left(\frac{1}{4}q_1^2 + 4q_2^2\right)^{-3/2}, \\
J_{12} &= q_2\left(8q_2^2 - q_1^2\right)\left(\frac{1}{4}q_1^2 + 4q_2^2\right)^{-3/2}, \\
J_{21} &= -6q_2^3\left(\frac{1}{4}q_1^2 + 4q_2^2\right)^{-3/2}, \\
J_{22} &= -\frac{3}{8}q_1^3\left(\frac{1}{4}q_1^2 + 4q_2^2\right)^{-3/2}.
\end{align*}

根据式(\ref{eq:9.2.14})，机器人机械臂方程由下式给出
\begin{equation}\tag{12}
\tau = M\ddot{q} + G + F(\dot{q}) + J^T(q)f
\end{equation}
其中$\tau$、$M$、$q$、$G$、$f$和$F(\dot{q})$如示例9.2.1所定义。重要的是要注意法向力（即$f_1$）和切向力（即$f_2$）是按(1)给出的任务空间坐标系的方向绘制的（见图9.2.6）。

%-------------------------------------------
\subsection{$N$连杆机械臂的刚度控制}
%-------------------------------------------

既然我们已经得到了包含环境相互作用力的机器人机械臂动力学形式，就可以建立$n$连杆机器人机械臂的刚度控制器。如前所述，施加在环境上的力定义为
\begin{equation}\label{eq:9.2.17}
f = K_e(x - x_e)
\end{equation}
其中$K_e$是一个$n \times n$对角、半正定常数矩阵，用于表示环境刚度，$x_e$是任务空间中测量的$n \times 1$向量，用于表示环境的静态位置。注意，如果机械臂在特定任务空间方向不受约束，则矩阵$K_e$的相应对角元素被假定为零。此外，在刚度控制公式中，环境表面摩擦通常被忽略。

多维刚度控制器是PD型控制器
\begin{equation}\label{eq:9.2.18}
\tau = J^T(q)\bigl(-K_v \dot{x} + K_p \tilde{x}\bigr) + G(q) + F(\dot{q})
\end{equation}
其中$K_v$和$K_p$是$n \times n$对角、常数、正定矩阵，任务空间跟踪误差定义为
\[
\tilde{x} = x_d - x
\]
如前所述，$x_d$用于表示我们希望将机器人机械臂移动到的期望恒定末端执行器位置；然而，$x_d$现在是一个$n \times 1$向量。将式(\ref{eq:9.2.17})和式(\ref{eq:9.2.18})代入式(\ref{eq:9.2.14})得到闭环动力学
\begin{equation}\label{eq:9.2.19}
M(q)\ddot{q} + V_m(q, \dot{q})\dot{q} = J^T(q)(-K_v\dot{x} + K_p\tilde{x} - K_e(x - x_e))
\end{equation}

为了分析式(\ref{eq:9.2.19})给出的系统的稳定性，我们利用类李雅普诺夫函数
\begin{equation}\label{eq:9.2.20}
V = \frac{1}{2}\dot{q}^T M(q) \dot{q} + \frac{1}{2}\tilde{x}^T K_p \tilde{x} + \frac{1}{2}(x - x_e)^T K_e (x - x_e)
\end{equation}
其中$\tilde{x} = x_d - x$是任务空间跟踪误差。将式(\ref{eq:9.2.20})关于时间微分并利用式(\ref{eq:9.2.10})得到
\begin{equation}\label{eq:9.2.21}
\dot{V} = \frac{1}{2}\dot{q}^T \dot{M}(q) \dot{q} - \dot{q}^T M(q) \ddot{q} - \dot{q}^T J^T(q) K_p \tilde{x} + \dot{q}^T J^T(q) K_e(x - x_e)
\end{equation}
注意，在式(\ref{eq:9.2.21})中，我们使用了$x_e$和$x_d$是常数的事实，并且标量函数或对角矩阵的转置分别等于该函数或矩阵本身。将式(\ref{eq:9.2.19})代入式(\ref{eq:9.2.21})并利用式(\ref{eq:9.2.10})得到
\begin{equation}\label{eq:9.2.22}
\dot{V} = \dot{q}^T\left(\frac{1}{2}\dot{M}(q) - V_m(q, \dot{q})\right)\dot{q} - \dot{q}^T J^T(q) K_v J(q) \dot{q}
\end{equation}
对式(\ref{eq:9.2.22})施用斜对称性质（见第2章）可得
\begin{equation}\label{eq:9.2.23}
\dot{V} = -\dot{q}^T J^T(q) K_v J(q) \dot{q}
\end{equation}
该$\dot{V}$为非正。由于矩阵$J(q)$以及因此$J^T(q) K_v J(q)$是非奇异的，$\dot{V}$只能在$\dot{q} = 0$因而$\ddot{q} = 0$的轨迹上恒为零（见第1章的拉萨尔定理）。将$\dot{q} = 0$与$\ddot{q} = 0$代入式(\ref{eq:9.2.19})并利用式(\ref{eq:9.2.10})得到
\begin{equation}\label{eq:9.2.24}
\lim_{t \to \infty} \bigl[K_p \tilde{x} - K_e(x - x_e)\bigr] = 0
\end{equation}
或等价地，
\begin{equation}\label{eq:9.2.25}
\lim_{t \to \infty} x_i = (K_{pi} + K_{ei})^{-1} (K_{pi} x_{di} + K_{ei} x_{ei})
\end{equation}
其中下标$i$用于表示向量$x$、$x_d$、$x_e$的第$i$个分量以及矩阵$K_p$和$K_e$的第$i$个对角元素。

上面的稳定性分析可以解释为意味着当任务空间坐标由式(\ref{eq:9.2.25})在稳态给出时，机器人机械臂将停止移动。也就是说，末端执行器的最终位置由式(\ref{eq:9.2.25})给出，这在单自由度情况下等价于由式(\ref{eq:9.2.6})给出。为了获得施加在环境上的稳态力的第$i$个分量，我们将式(\ref{eq:9.2.25})代入式(\ref{eq:9.2.17})的第$i$个分量，得到
\begin{equation}\label{eq:9.2.26}
\lim_{t \to \infty} f_i = K_{ei} (K_{pi} + K_{ei})^{-1} K_{pi} (x_{di} - x_{ei})
\end{equation}

因此末端执行器施加在环境上的稳态力由式(\ref{eq:9.2.26})给出，这在单自由度情况下等价于由式(\ref{eq:9.2.7})给出。如在单自由度情况下一样，我们假设对于要力控制的任务空间方向，$K_{ei}$远大于$K_{pi}$。也就是说，式(\ref{eq:9.2.26})中的稳态力可以近似为
\begin{equation}\label{eq:9.2.27}
\lim_{t \to \infty} f_i \cong K_{pi}(x_{di} - x_{ei})
\end{equation}
因此，$K_{pi}$可以被解释为在这些任务空间方向上指定机械臂的刚度。

如果机械臂在任务空间方向上不受约束，则相应的刚度常数$K_{ei}$等于零。将$K_{ei} = 0$代入式(\ref{eq:9.2.25})得到
\begin{equation}\label{eq:9.2.28}
\lim_{t \to \infty} x_i = x_{di}
\end{equation}
这意味着对于不受约束的任务空间方向，我们获得了设定点控制；因此，在稳态下达到期望的位置设定点。刚度控制器以及相应的稳定性结果都在表9.2.1中总结。

\textbf{表9.2.1}\quad 刚度控制器

\begin{center}
\small
\begin{tabular}{|l|p{11cm}|}
\hline
\textbf{扭矩控制器} & $\displaystyle \tau = J^{\mathsf T}(q)\bigl(-K_v \dot{x} + K_p \tilde{x}\bigr) + G(q) + F(\dot{q})$，其中$J(q)$为任务空间雅可比矩阵，$J^{\mathsf T}(q)$为其转置。 \\
\hline
\textbf{稳定性} &
\setlength{\parskip}{6pt}
\textbf{(i)}\ 不受约束方向（设定点位置控制）：$\displaystyle \lim_{t\to\infty} x_i(t) = x_{di}$。\par
\textbf{(ii)}\ 受约束方向（稳态力近似）：$\displaystyle \lim_{t\to\infty} f_i(t) \cong K_{pi}(x_{di} - x_{ei})$。 \\
\hline
\textbf{注释} & 控制增益$K_{pi}$用于调节机械臂刚度。 \\
\hline
\end{tabular}
\end{center}

\textbf{示例 9.2--3：笛卡尔机械臂的刚度控制器}

我们希望为图9.2.5中给出的机器人机械臂系统设计并模拟一个刚度控制器。控制目标是将末端执行器移动到期望的最终切向位置$v_d = 3\ \mathrm{m}$，同时施加$f_{d1} = 2\ \mathrm{N}$的最终期望法向力。我们忽略表面摩擦（即$f_2$）和关节摩擦；连杆质量取为单位质量；并设法向力满足
\begin{equation}\tag{1}
f_1 = k_e(u - u_e)
\end{equation}
其中$u = \frac{1}{\sqrt{2}}(q_1 - q_2)$，环境刚度$k_e = 1000\ \mathrm{N/m}$，环境平衡位置
\begin{equation*}
u_e = \frac{3}{\sqrt{2}}\ \mathrm{m}.
\end{equation*}
初始末端执行器任务坐标取为
\begin{equation}\tag{2}
u(0) = \frac{3}{\sqrt{2}}\ \mathrm{m}, \quad v(0) = 5\ \mathrm{m}, \quad v = \frac{1}{\sqrt{2}}(q_1 + q_2).
\end{equation}

刚度控制器（忽略摩擦项）为
\begin{equation}\tag{3}
\tau = J^{\mathsf T}(q)\bigl(-K_v\dot{x} + K_p\tilde{x}\bigr) + G(q)
\end{equation}
其中$\tau$、$J(q)$、$G(q)$与$x=[u,\ v]^{\mathsf T}$如示例9.2.1所定义，
\begin{equation*}
\tilde{x} = \begin{bmatrix} u_d - u \\ v_d - v \end{bmatrix},
\end{equation*}
$u_d$为待定的期望法向位置，增益取$K_v = k_v I$、$K_p = k_p I$，且$k_v = k_p = 10$，从而满足刚度控制要求$k_p \ll k_e$。由式(\ref{eq:9.2.27})的近似，用法向期望力确定$u_d$时采用
\begin{equation}\tag{4}
f_{d1} = k_p(u_d - u_e),
\end{equation}
代入$f_{d1} = 2$、$k_p = 10$与$u_e = 3/\sqrt{2}$得
\begin{align*}
2 &= 10\left(u_d - \frac{3}{\sqrt{2}}\right), \\
u_d &= 0.2 + \frac{3}{\sqrt{2}} \approx 2.321\ \mathrm{m}.
\end{align*}

图9.2.7所示机械臂系统（图9.2.5）按式(3)仿真表明，期望切向位置与法向力约在4秒内达到。

%===========================================
\section{混合位置/力控制}\label{sec:9.3}
%===========================================

第\ref{sec:9.2}节给出的刚度控制器的一个主要缺点是它只能用于设定点控制；换句话说，期望的末端执行器机械臂位置和施加在环境上的期望力必须是常数。在许多机器人应用中，例如磨削，末端执行器必须跟踪沿物体表面的期望位置轨迹，同时跟踪施加到物体表面的期望力轨迹。在这种类型的应用中，刚度控制器将不能充分执行；因此，必须使用另一种控制方法。

所谓的混合位置/力控制器[Chae et al. 1988]和[Raibert and Craig 1981]可用于同时跟踪位置和力轨迹。混合位置/力控制器的基本概念是通过任务空间公式将位置和力控制问题解耦为子任务。正如我们所见，任务空间公式在确定哪些方向应该受力或位置控制方面很有价值。也就是说，位置和力控制子任务很容易从任务空间公式中确定。在确定控制子任务后，可以分别开发位置控制器和力控制器。

%-------------------------------------------
\subsection{笛卡尔二连杆臂的混合位置/力控制}
%-------------------------------------------

为了说明这种混合位置/力控制的概念，考虑图9.3.1给出的机器人机械臂系统。对于这个应用，沿表面的位置和施加在表面上的法向力都应该被控制；因此，必须确定哪些变量应该受力控制，哪些应该受位置控制。

根据第\ref{sec:9.2}节给出的任务空间概念，图9.3.1给出的机械臂系统的任务空间公式为
\begin{equation}\label{eq:9.3.1}
x = \begin{bmatrix} u \\ v \end{bmatrix} = h(q) = \begin{bmatrix} q_1 \\ q_2 \end{bmatrix}
\end{equation}
任务空间雅可比矩阵由下式给出
\[
J = \begin{bmatrix} 1 & 0 \\ 0 & 1 \end{bmatrix}
\]

如式(\ref{eq:9.3.1})所示，对于这个问题的任务空间和关节空间是等价的；因此，我们将在整个问题中将关节变量称为任务空间变量。

为了为机械臂系统设计位置/力控制器，我们必须首先确定式(\ref{eq:9.3.1})给出的任务空间公式的动力学方程。使用这个任务空间公式并忽略关节摩擦，可以证明机械臂动力学为
\begin{equation}\label{eq:9.3.2}
\tau = M\ddot{q} + G + f
\end{equation}
其中$\tau$、$M$、$G$和$f$如示例9.2.1所定义。式(\ref{eq:9.3.2})表示的矩阵形式中的两个动力学方程是
\begin{equation}\label{eq:9.3.3}
\tau_1 = m_1\ddot{q}_1 + f_1
\end{equation}
和
\begin{equation}\label{eq:9.3.4}
\tau_2 = (m_1 + m_2)\ddot{q}_2 + (m_1 + m_2)g + f_2
\end{equation}

在制定混合位置/力控制器时，我们为式(\ref{eq:9.3.3})和式(\ref{eq:9.3.4})给出的动力学设计了单独的控制器。如图9.3.1所示，沿任务空间方向$q_2$的位置应该受位置控制；因此，我们应该使用式(\ref{eq:9.3.4})来设计位置控制器。[这很明显，因为式(\ref{eq:9.3.3})给出的动力学不包含任务空间变量$q_2$。]

因为我们正在设计一个位置控制器来跟踪期望轨迹，我们将定义``切向空间''跟踪误差为
\begin{equation}\label{eq:9.3.5}
\tilde{x} = q_{d2} - q_2
\end{equation}
其中$q_{d2}$表示沿表面或切向表面的期望位置轨迹。位置控制器将是计算力矩控制器（见第3章）
\begin{equation}\label{eq:9.3.6}
\tau_2 = (m_1 + m_2)a_T + (m_1 + m_2)g + f_2
\end{equation}
其中
\begin{equation}\label{eq:9.3.7}
a_T = \ddot{q}_{d2} + k_{Tv}\dot{\tilde{x}} + k_{Tp}\tilde{x}
\end{equation}
$k_{Tv}$和$k_{Tp}$是正控制增益。将式(\ref{eq:9.3.6})代入式(\ref{eq:9.3.4})得到位置跟踪误差系统
\begin{equation}\label{eq:9.3.8}
\ddot{\tilde{x}} + k_{Tv}\dot{\tilde{x}} + k_{Tp}\tilde{x} = 0
\end{equation}

利用$k_{Tv}$和$k_{Tp}$是正数的事实，我们可以对式(\ref{eq:9.3.8})应用标准线性控制结果，得到
\[
\lim_{t \to \infty} \tilde{x} = 0
\]
因此，式(\ref{eq:9.3.6})给出的控制器保证了渐近位置跟踪。注意，位置控制器需要测量关节位置、关节速度和表面摩擦力；因此，从实现的角度来看，位置控制器中需要力测量。

式(\ref{eq:9.3.6})给出的位置控制器将确保沿环境表面的良好位置跟踪；然而，我们还希望控制施加在环境上的力。对于图9.3.1给出的机械臂系统，任务空间方向垂直于表面的是$q_1$；因此，我们将假设在这个方向上环境可以被建模为弹簧。具体来说，施加在环境上的法向力$f_1$由下式给出
\begin{equation}\label{eq:9.3.9}
f_1 = k_e(q_1 - q_e)
\end{equation}
其中$k_e$表示环境刚度，$q_e = 3$。将式(\ref{eq:9.3.9})关于时间的二阶导数得到表达式
\begin{equation}\label{eq:9.3.10}
\ddot{q}_1 = \frac{1}{k_e}\ddot{f}_1
\end{equation}
其中法向任务空间加速度用设法向力的二阶导数来书写。将式(\ref{eq:9.3.10})代入式(\ref{eq:9.3.3})得到力动力学方程
\begin{equation}\label{eq:9.3.11}
\tau_1 = \frac{m_1}{k_e}\ddot{f}_1 + f_1
\end{equation}

我们现在可以使用式(\ref{eq:9.3.11})给出的力动力学方程来设计一个力控制器来跟踪期望的力轨迹。首先，定义力跟踪误差为
\begin{equation}\label{eq:9.3.12}
\tilde{f} = f_{d1} - f_1
\end{equation}
其中$f_{d1}$表示要施加在环境上的期望法向力。与位置控制器类似，力控制器将是计算力矩控制器
\begin{equation}\label{eq:9.3.13}
\tau_1 = \frac{m_1}{k_e}a_N + f_1
\end{equation}
其中
\begin{equation}\label{eq:9.3.14}
a_N = \ddot{f}_{d1} + k_{Nv}\dot{\tilde{f}} + k_{Np}\tilde{f}
\end{equation}
$k_{Nv}$和$k_{Np}$是正控制增益。将式(\ref{eq:9.3.13})代入式(\ref{eq:9.3.11})得到力跟踪误差系统
\begin{equation}\label{eq:9.3.15}
\ddot{\tilde{f}} + k_{Nv}\dot{\tilde{f}} + k_{Np}\tilde{f} = 0
\end{equation}
利用式(\ref{eq:9.3.15})中$k_{Nv}$和$k_{Np}$是正数的事实，得到
\[
\lim_{t \to \infty} \tilde{f} = 0
\]
因此，式(\ref{eq:9.3.13})给出的控制器保证了渐近力跟踪。重要的是要认识到力控制器需要测量法向力及其导数。因为力导数通常不用于测量，所以由式(\ref{eq:9.3.9})制造，即
\begin{equation}\label{eq:9.3.16}
\dot{f}_1 = k_e\dot{q}_1
\end{equation}
因此，环境刚度和法向任务空间速度用于模拟力的导数。

式(\ref{eq:9.3.6})给出的位置控制器和式(\ref{eq:9.3.13})给出的力控制器构成了完整的混合位置/力控制器。这将在下一小节中推广到$N$连杆机械臂。

%-------------------------------------------
\subsection{$N$连杆机械臂的混合位置/力控制}
%-------------------------------------------

上一节给出的混合位置/力控制器可以通过使用任务空间公式概念轻松扩展到多自由度情况。具体来说，可以开发一个反馈线性化控制，使机器人机械臂方程全局线性化，然后开发线性控制器来跟踪期望的力和位置轨迹。

首先，控制设计者选择一个任务空间公式
\begin{equation}\label{eq:9.3.17}
x = h(q)
\end{equation}
使得法向和切向表面运动如第\ref{sec:9.2}节所讨论的那样被分解。式(\ref{eq:9.2.14})给出的机器人动力学然后用任务空间加速度表示，通过对式(\ref{eq:9.3.17})关于时间求二阶导数得到
\begin{equation}\label{eq:9.3.18}
\ddot{x} = J(q)\ddot{q} + \dot{J}(q)\dot{q}
\end{equation}
其中$J(q)$是式(\ref{eq:9.2.11})中定义的任务空间雅可比矩阵。求解式(\ref{eq:9.3.18})中的$\ddot{q}$得到
\begin{equation}\label{eq:9.3.19}
\ddot{q} = J^{-1}(q)(\ddot{x} - \dot{J}(q)\dot{q})
\end{equation}

将式(\ref{eq:9.3.19})代入式(\ref{eq:9.2.14})并重新排列，我们可以证明机器人动力学可以写成形式
\begin{equation}\label{eq:9.3.20}
\tau = M(q)J^{-1}(q)(\ddot{x} - \dot{J}(q)\dot{q}) + V_m(q, \dot{q})\dot{q} + G(q) + F(\dot{q}) + J^T(q)f
\end{equation}

式(\ref{eq:9.3.20})给出的动力学对应的反馈线性化控制由下式给出
\begin{equation}\label{eq:9.3.21}
\tau = M(q)J^{-1}(q)(a - \dot{J}(q)\dot{q}) + V_m(q, \dot{q})\dot{q} + G(q) + F(\dot{q}) + J^T(q)f
\end{equation}
其中$a$是一个$n \times 1$向量，用于表示线性位置和力控制策略，稍后将讨论。将式(\ref{eq:9.3.21})代入式(\ref{eq:9.3.20})后，我们有
\begin{equation}\label{eq:9.3.22}
\ddot{x} = a
\end{equation}

从式(\ref{eq:9.3.22})中，我们可以看到任务空间运动已经被全局线性化和解耦；因此，我们可以为每个任务空间自由度独立设计位置和力控制器。具体来说，可以为表示切向运动的任务空间变量设计线性位置控制器；此外，可以为表示法向力的任务空间变量设计线性力控制器。

因为式(\ref{eq:9.3.20})给出的动力学在任务空间中已经被解耦，我们将$x$的切向空间分量定义为$x_{Ti}$，其中下标$T$用于表示切向空间，下标$i$用于表示$x_T$的第$i$个分量。根据这个符号，式(\ref{eq:9.3.22})的切向空间分量由下式给出
\begin{equation}\label{eq:9.3.23}
\ddot{x}_{Ti} = a_{Ti}
\end{equation}
其中$a_{Ti}$是第$i$个线性切向空间位置控制器。为了反馈控制的目的，我们定义切向空间跟踪误差为
\begin{equation}\label{eq:9.3.24}
\tilde{x}_{Ti} = x_{Tdi} - x_{Ti}
\end{equation}
其中$x_{Tdi}$表示沿环境表面切向的第$i$个期望位置轨迹。如上一节所示，相应的线性控制器由下式给出
\begin{equation}\label{eq:9.3.25}
a_{Ti} = \ddot{x}_{Tdi} + k_{Tvi}\dot{\tilde{x}}_{Ti} + k_{Tpi}\tilde{x}_{Ti}
\end{equation}
$k_{Tvi}$和$k_{Tpi}$是第$i$个正控制增益。将式(\ref{eq:9.3.25})代入式(\ref{eq:9.3.23})得到位置跟踪误差系统
\begin{equation}\label{eq:9.3.26}
\ddot{\tilde{x}}_{Ti} + k_{Tvi}\dot{\tilde{x}}_{Ti} + k_{Tpi}\tilde{x}_{Ti} = 0
\end{equation}
利用式(\ref{eq:9.3.26})中$k_{Tvi}$和$k_{Tpi}$是正数的事实，得到
\[
\lim_{t \to \infty} \tilde{x}_{Ti} = 0
\]
因此，保证了渐近位置跟踪。

为了力控制的目的，我们将$x$的法向空间分量定义为$x_{Nj}$，其中下标$N$用于表示法向空间，下标$j$用于表示$x_N$的第$j$个分量。根据这个符号，式(\ref{eq:9.3.22})的法向空间分量由下式给出
\begin{equation}\label{eq:9.3.27}
\ddot{x}_{Nj} = \mathbf{a}_{Nj}
\end{equation}
其中$\mathbf{a}_{Nj}$是第$j$个线性法向空间力控制器。如上一节所示，我们假设环境可以被建模为弹簧。具体来说，施加在环境上的法向力$f_{Nj}$由下式给出
\begin{equation}\label{eq:9.3.28}
f_{Nj} = k_{ej}(x_{Nj} - x_{ej})
\end{equation}
其中$k_{ej}$是环境刚度的第$j$个分量，$x_{ej}$用于表示法向空间$x_{Nj}$方向上的环境静态位置。

如对上一小节中的单自由度机器人所做的那样，在开发力控制器之前，我们必须制定力动力学。对式(\ref{eq:9.3.28})关于时间求二阶导数得到表达式
\begin{equation}\label{eq:9.3.29}
\ddot{f}_{Nj} = k_{ej}\ddot{x}_{Nj}
\end{equation}
其中法向任务空间加速度用第$j$个法向力的二阶导数来书写。将式(\ref{eq:9.3.29})代入式(\ref{eq:9.3.27})得到力动力学
\begin{equation}\label{eq:9.3.30}
\frac{1}{k_{ej}}\ddot{f}_{Nj} = \mathbf{a}_{Nj}
\end{equation}

为了反馈控制的目的，我们定义力跟踪误差为
\begin{equation}\label{eq:9.3.31}
\tilde{f}_{Nj} = f_{Ndj} - f_{Nj}
\end{equation}
其中$f_{Ndj}$表示要垂直施加在环境上的期望力的第$j$个分量。如上一节所示，相应的线性控制器由下式给出
\begin{equation}\label{eq:9.3.32}
\mathbf{a}_{Nj} = \frac{1}{k_{ej}}(\ddot{f}_{Ndj} + k_{Nvj}\dot{\tilde{f}}_{Nj} + k_{Npj}\tilde{f}_{Nj})
\end{equation}
$k_{Nvj}$和$k_{Npj}$是第$j$个正控制增益。将式(\ref{eq:9.3.32})代入式(\ref{eq:9.3.30})得到力跟踪误差系统
\begin{equation}\label{eq:9.3.33}
\ddot{\tilde{f}}_{Nj} + k_{Nvj}\dot{\tilde{f}}_{Nj} + k_{Npj}\tilde{f}_{Nj} = 0
\end{equation}

利用式(\ref{eq:9.3.33})中$k_{Nvj}$和$k_{Npj}$是正数的事实，得到
\[
\lim_{t \to \infty} \tilde{f}_{Nj} = 0
\]
因此，保证了渐近力跟踪。

混合位置/力控制器以及相应的稳定性结果都在表9.3.1中总结。

\textbf{表9.3.1}\quad 混合位置/力控制器

\begin{center}
\small
\begin{tabular}{|l|p{11.2cm}|}
\hline
\textbf{扭矩控制器} & \begin{minipage}[t]{\linewidth}\raggedright\setlength{\lineskip}{6pt}
$\displaystyle \tau = M(q)J^{-1}(q)\bigl(a - \dot{J}(q)\dot{q}\bigr) + V_m(q,\dot{q})\dot{q} + G(q) + F(\dot{q}) + J^{\mathsf T}(q)f$，其中$J(q)$为任务空间雅可比；$n\times 1$向量$a$表示线性位置/力控制策略，其分量由下列位置控制与力控制式给出。\par
\end{minipage} \\
\hline
\textbf{位置控制}（非约束方向$i$） & $\displaystyle a_{Ti} = \ddot{x}_{Tdi} + k_{Tvi}\dot{\tilde{x}}_{Ti} + k_{Tpi}\tilde{x}_{Ti}$ \\
\hline
\textbf{力控制}（约束方向$j$） & $\displaystyle a_{Nj} = \frac{1}{k_{ej}}\bigl(\ddot{f}_{Ndj} + k_{Nvj}\dot{\tilde{f}}_{Nj} + k_{Npj}\tilde{f}_{Nj}\bigr)$ \\
\hline
\textbf{稳定性} &
\setlength{\parskip}{6pt}
\textbf{(i)}\ 非约束方向（位置跟踪）：$\displaystyle \lim_{t\to\infty} x_{Ti}(t) = x_{Tdi}(t)$。\par
\textbf{(ii)}\ 约束方向（力跟踪）：$\displaystyle \lim_{t\to\infty} f_{Nj}(t) = f_{Ndj}(t)$。 \\
\hline
\textbf{注释} & 环境建模为弹簧。 \\
\hline
\end{tabular}
\end{center}

\textbf{示例 9.3--1：倾斜表面上笛卡尔机械臂的混合位置/力控制}

我们希望为图9.2.5中所示机器人机械臂系统设计并仿真混合位置/力控制器，使末端执行器沿表面跟踪期望切向位置轨迹$v_d = \sin t\ \mathrm{(m)}$，同时跟踪期望法向力轨迹$f_{Nd1} = 1 - \mathrm{e}^{-t}\ \mathrm{(N)}$。忽略关节摩擦；连杆质量取为单位质量；设法向力满足
\begin{equation}\label{eq:9.3.38}
f_1 = k_e(u - u_e)
\end{equation}
其中$u = \frac{1}{\sqrt{2}}(q_1 - q_2)$，$k_e = 1000\ \mathrm{N/m}$，环境平衡位置
\begin{equation*}
u_e = \frac{3}{\sqrt{2}}\ \mathrm{m}.
\end{equation*}
初始任务坐标取为
\begin{equation}\label{eq:9.3.39}
v(0) = 0, \quad u(0) = \frac{3}{\sqrt{2}}\ \mathrm{m}, \quad v = \frac{1}{\sqrt{2}}(q_1 + q_2).
\end{equation}

控制器结构见表9.3.1。与示例9.2.1一致，当忽略科氏/离心项及关节摩擦时，扭矩律可写为
\begin{equation}\label{eq:9.3.40}
\tau = M J^{-1} a + G + J^{\mathsf T} f
\end{equation}
其中$a$为$2\times 1$线性加速度指令向量，$M$、$J$、$G$如示例9.2.1；闭环任务空间方程为
\begin{equation}\label{eq:9.3.41}
\ddot{x} = \begin{bmatrix} \ddot{u} \\ \ddot{v} \end{bmatrix} = a.
\end{equation}
此处$u$为法向坐标（约束方向、力控），$v$为切向坐标（非约束方向、位置控），故
\begin{equation}\label{eq:9.3.42}
\ddot{x} = \begin{bmatrix} \ddot{x}_{N1} \\ \ddot{x}_{T1} \end{bmatrix} = \begin{bmatrix} a_{N1} \\ a_{T1} \end{bmatrix} = a.
\end{equation}
其中$x_{T1}=v$，$x_{N1}=u$，$\tilde{x}_{T1}=x_{Td1}-x_{T1}$，$\tilde{f}_{N1}=f_{Nd1}-f_1$。
由表9.3.1，切向与法向线性控制律分别为
\begin{equation}\label{eq:9.3.43}
a_{T1} = \ddot{x}_{Td1} + k_{Tv1}\dot{\tilde{x}}_{T1} + k_{Tp1}\tilde{x}_{T1}
\end{equation}
\begin{equation}\label{eq:9.3.44}
a_{N1} = \frac{1}{k_{e1}}\bigl(\ddot{f}_{Nd1} + k_{Nv1}\dot{\tilde{f}}_{N1} + k_{Np1}\tilde{f}_{N1}\bigr)
\end{equation}
仿真中取期望轨迹与环境参数
\begin{equation}\label{eq:9.3.45}
x_{Td1} = \sin t, \quad f_{Nd1} = 1 - \mathrm{e}^{-t}, \quad k_{e1} = 1000,
\end{equation}
\begin{equation*}
k_{Nv1} = k_{Np1} = k_{Tv1} = k_{Tp1} = 10.
\end{equation*}
图9.3.2所示结果表明，位置与力跟踪误差约在$4\ \mathrm{s}$内趋于零。

如前所述，混合位置/力控制器需要测量力变量$f$。如前所述，可以使用变换来获得任何特定力应用的$f$测量值。也就是说，任务空间力与传感器力的关系为
\begin{equation}\label{eq:9.3.34}
f = J_s^T(q) f_s
\end{equation}
其中$J_s(q)$是一个$n \times n$传感器雅可比矩阵，$f_s$是一个$n \times 1$传感器力向量。使用式(\ref{eq:9.3.34})，任务空间力由下式给出
\begin{equation}\label{eq:9.3.35}
f = J_s^T(q) f_s
\end{equation}
正如我们前面提到的，力控制律需要测量任务空间力导数；然而，这个信号通常不可用。通常，刚度方程(\ref{eq:9.3.28})用于获得第$j$个任务空间力导数
\begin{equation}\label{eq:9.3.36}
\dot{f}_{Nj} = k_{ej}\dot{x}_{Nj}
\end{equation}
它是用可测量的第$j$个任务空间法向速度表示的。

考虑到这些实现问题，整体混合位置/力控制策略如图9.3.3所示。框图中的前馈项用于表示控制律中的项
\begin{equation}\label{eq:9.3.37}
-M(q)J^{-1}(q)\dot{J}(q)\dot{q} + V_m(q,\dot{q})\dot{q} + G(q) + F(\dot{q}) + J^{\mathsf T}(q)f
\end{equation}
（与各方向的加速度指令分量$a_{Ti}$、$a_{Nj}$——参见式(\ref{eq:9.3.25})与式(\ref{eq:9.3.32})——共同出现于完整的扭矩合成中。）

%===========================================
\section{混合阻抗控制}\label{sec:9.4}
%===========================================

阻抗控制基于这样的概念：控制器应用于调节机器人机械臂运动与施加在环境上的力之间的动态行为[Hogan 1987]，而不是分别考虑运动和力控制问题。使用这个概念，控制设计者指定机械臂运动和施加在环境上的力之间的期望动态行为。这种期望行为有时被称为目标阻抗，因为它用于表示运动和力之间的欧姆定律类型关系。

%-------------------------------------------
\subsection{环境建模}
%-------------------------------------------

正如[Hogan 1987]中所指出的，环境模型对任何力控制策略都是至关重要的。在前面讨论的力控制策略中，环境简单地被建模为弹簧；然而，可以想象，简单的弹簧模型可能不足以描述所有类型的环境。为了对多种类型的环境进行分类，我们使用线性传递函数关系
\begin{equation}\label{eq:9.4.1}
f(s) = Z_e(s)\dot{x}(s)
\end{equation}
其中变量$s$是拉普拉斯变换变量，$f$表示施加在环境上的力，$\dot{x}$表示机械臂在环境接触点的速度，$Z_e(s)$表示环境阻抗。目前，所有量都假定为标量函数；然而，在本节末尾，我们将阻抗控制推广到多维情况。

量$Z_e(s)$被称为阻抗，因为式(\ref{eq:9.4.1})表示运动和力之间的欧姆定律类型关系。如在电路理论中，环境阻抗可以分为不同的类别。为了进一步讨论阻抗控制，我们现在给出三种常用的用于分类环境阻抗的类别。

\begin{definition}\label{def:9.4.1}
当且仅当$|Z(0)| = 0$时，阻抗是惯性的。
\end{definition}

惯性环境的说明由图9.4.1a给出。该图描绘了机器人机械臂以速度$\dot{x}$移动质量为$h$的载荷。相应的相互作用力由下式给出
\[
f = h\ddot{x}
\]
因此，利用式(\ref{eq:9.4.1})得到惯性环境阻抗为
\begin{equation}\label{eq:9.4.2}
Z_e(s) = hs
\end{equation}
我们可以通过应用定义\ref{def:9.4.1}轻松验证这个阻抗确实是惯性的。

\begin{definition}\label{def:9.4.2}
当且仅当$|Z(0)| = c$时，阻抗是电阻性的，其中$0 < c < \infty$。
\end{definition}

电阻性环境的说明由图9.4.1b给出。该图描绘了机器人机械臂以速度$\dot{x}$穿过液体介质。液体介质被假定为具有阻尼系数$b$。相应的相互作用力由下式给出
\[
f = b\dot{x}
\]
因此，利用式(\ref{eq:9.4.1})得到电阻性环境阻抗为
\begin{equation}\label{eq:9.4.3}
Z_e(s) = b
\end{equation}
我们可以通过应用定义\ref{def:9.4.2}轻松验证这个阻抗确实是电阻性的。

\begin{definition}\label{def:9.4.3}
当且仅当$|Z(0)| = \infty$时，阻抗是电容性的。
\end{definition}

电容性环境的说明由图9.4.1c给出。该图描绘了机器人机械臂以速度$\dot{x}$推动质量为$h$的物体。物体被假定为具有阻尼系数$b$和弹簧常数$k$。相应的相互作用力由下式给出
\[
f = h\ddot{x} + b\dot{x} + kx
\]
因此，利用式(\ref{eq:9.4.1})得到电容性环境阻抗为
\begin{equation}\label{eq:9.4.4}
Z_e(s) = hs + b + \frac{k}{s}
\end{equation}
我们可以通过应用定义\ref{def:9.4.3}轻松验证这个阻抗确实是电容性的。

%-------------------------------------------
\subsection{位置和力控制模型}
%-------------------------------------------

正如我们在前一节中看到的，环境可以被建模为由式(\ref{eq:9.4.1})中的力/速度关系定义的阻抗。对于我们的阻抗控制公式，我们将假设环境阻抗是惯性的、电阻性的或电容性的。问题就变成了：如何为给定的环境阻抗设计控制器？解决方案是通过制定机械臂阻抗模型$Z_m(s)$[Anderson and Spong 1988]来获得的。换句话说，在对环境建模之后选择机械臂阻抗（或目标阻抗）。选择机械臂阻抗的标准与机械臂的动态性能有关。也就是说，选择机械臂阻抗使得对阶跃输入（可能是力或速度命令）的稳态误差为零。正如我们将展示的，如果机械臂阻抗是环境阻抗的对偶，就可以实现这个性能标准。

在完全阐述对偶概念之前，必须制定位置和力控制的模型。对于位置控制[Anderson and Spong 1988]，力与速度之间的关系由下式建模
\begin{equation}\label{eq:9.4.5}
f = Z_m(s)(\dot{x}_d - \dot{x})
\end{equation}
其中$\dot{x}_d$表示机械臂在环境接触点的输入速度，$Z_m(s)$表示机械臂阻抗。

如我们随后将展示的，机械臂阻抗$Z_m(s)$被选择为通过对阶跃输入利用$x$和$\dot{x}_d$之间的动态关系来``归零''稳态。为了确定$x$和$\dot{x}_d$之间的动态关系，我们将式(\ref{eq:9.4.1})与式(\ref{eq:9.4.5})组合，得到图9.4.2给出的位置控制框图。我们可以使用这个框图来说明对偶概念。具体来说，我们检查稳态速度误差
\begin{equation}\label{eq:9.4.6}
E_{ss} = \lim_{s \to 0} s(\dot{x}_d - \dot{x})
\end{equation}
其中$\dot{x}_d(s) = 1/s$对于阶跃速度输入。利用图9.4.2，我们可以很容易地证明式(\ref{eq:9.4.6})可以简化为
\begin{equation}\label{eq:9.4.7}
E_{ss} = \lim_{s \to 0} \frac{Z_e(s)}{Z_e(s) + Z_m(s)}
\end{equation}

对于式(\ref{eq:9.4.7})中的$E_{ss}$等于零，$Z_e(s)$必须是非电容性阻抗[即$Z_e(s)$必须是惯性或电阻性的]，$Z_m(s)$必须是非惯性阻抗。也就是说，如果惯性环境用非惯性机械臂阻抗进行位置控制，而电阻性环境用电容性机械臂阻抗进行位置控制，则对速度阶跃输入可以实现零稳态误差。上述术语``对偶''用于强调惯性环境阻抗可以用电容性机械臂阻抗进行位置控制。

从上面的发展来看，很明显电容性环境不能被位置控制并保持零稳态误差规范。然而，随后我们将展示电容性环境可以在保持零稳态误差规范的同时被力控制。关于力控制，力与速度之间的动态关系被建模为[Anderson and Spong 1988]
\begin{equation}\label{eq:9.4.8}
\dot{x} = Z_m^{-1}(s)(f_d - f)
\end{equation}
其中$f_d$用于表示在环境接触点施加的输入力。

如我们随后将展示的，机械臂阻抗$Z_m(s)$被选择为通过对阶跃输入利用$f$和$f_d$之间的动态关系来``归零''稳态。为了确定$f$和$f_d$之间的动态关系，我们将式(\ref{eq:9.4.1})与式(\ref{eq:9.4.8})组合，得到图9.4.3给出的力控制框图。我们可以使用这个框图来说明对偶概念。具体来说，我们检查稳态力误差
\begin{equation}\label{eq:9.4.9}
E_{ss} = \lim_{s \to 0} s(f_d - f)
\end{equation}
其中$f_d(s) = 1/s$对于阶跃力输入。利用图9.4.3，我们可以很容易地证明式(\ref{eq:9.4.9})可以简化为
\begin{equation}\label{eq:9.4.10}
E_{ss} = \lim_{s \to 0} \frac{Z_m(s)}{Z_e(s) + Z_m(s)}
\end{equation}

对于式(\ref{eq:9.4.10})中的$E_{ss}$等于零，$Z_e(s)$必须是非惯性阻抗，$Z_m(s)$必须是非电容性阻抗。也就是说，如果电容性环境用非电容性机械臂阻抗进行力控制，而电阻性环境用惯性机械臂阻抗进行力控制，则对力阶跃输入可以实现零稳态误差。术语``对偶''用于强调电容性环境阻抗可以用惯性机械臂阻抗进行力控制。

上面的讨论可以由以下对偶原理总结。

\textbf{对偶原理} 电容性环境用非电容性机械臂阻抗进行力控制，惯性环境用非惯性机械臂阻抗进行位置控制，电阻性环境用惯性机械臂阻抗进行力控制或用电容性机械臂阻抗进行位置控制。

\begin{quote}
\textit{注：读者可以通过}\href{https://frankjie09.github.io/Robot-Manipulator-Control-CN/figures/Impedance_Control.html}{\textbf{交互式可视化工具}}\textit{理解阻抗控制的概念，调节虚拟刚度、阻尼和质量参数，观察机器人对外力的顺应响应。}
\end{quote}

%-------------------------------------------
\subsection{阻抗控制公式}
%-------------------------------------------

既然我们已经展示了如何基于这个模型将环境和机械臂建模为阻抗，我们就基于这个模型开发一个``阻抗''控制器。为了利用阻抗控制方法，控制设计者选择一个任务空间公式
\begin{equation}\label{eq:9.4.11}
x = h(q)
\end{equation}
用于特定的位置/力控制应用。如第\ref{sec:9.3}节所示，我们可以证明力矩控制
\begin{equation}\label{eq:9.4.12}
\tau = M(q)J^{-1}(q)(a - \dot{J}(q)\dot{q}) + V_m(q, \dot{q})\dot{q} + G(q) + F(\dot{q}) + J^T(q)f
\end{equation}
得到线性方程组
\begin{equation}\label{eq:9.4.13}
\ddot{x} = a
\end{equation}
其中$a$是一个$n \times 1$向量，用于表示阻抗位置和力控制策略。

如式(\ref{eq:9.4.13})所阐述的，任务空间运动已经被全局线性化；因此，我们可以为每个任务空间自由度设计单独的位置和力控制器。也就是说，在每个由分量$x_k$表示的任务空间方向上，分配$x_k$与相应环境力$f_k$之间的环境阻抗关系。基于这种环境阻抗的分配，对偶原理用于确定$a$的相应元素应该是力控制器还是位置控制器。在做出这个决定后，然后使用式(\ref{eq:9.4.5})和式(\ref{eq:9.4.8})来获得$a$的特定位置和力控制分量。

作为将位置控制设计与力控制设计分离的手段，我们使用式(\ref{eq:9.4.13})来定义在任务空间方向上受位置控制的方程为
\begin{equation}\label{eq:9.4.14}
\ddot{x}_{pi} = a_{pi}
\end{equation}
其中下标$i$表示第$i$个受位置控制的任务空间变量，下标$p$表示位置控制。与受位置控制的任务空间方向相关的环境力用$f_{pi}$表示。

假设零初始条件，式(\ref{eq:9.4.14})的拉普拉斯变换可以写成
\begin{equation}\label{eq:9.4.15}
s\dot{x}_{pi} = a_{pi}
\end{equation}
从式(\ref{eq:9.4.5})给出的位置控制模型，其第$i$个分量满足
\[
f_{pi} = Z_{pmi}(s)\bigl(\dot{x}_{pdi} - \dot{x}_{pi}\bigr),
\]
故$\dot{x}_{pi} = \dot{x}_{pdi} - Z_{pmi}^{-1}(s)f_{pi}$。两端乘以$s$后，式(\ref{eq:9.4.15})的左侧亦可写成
\begin{equation}\label{eq:9.4.16}
s\dot{x}_{pi}(s) = s\left(\dot{x}_{pdi}(s) - Z_{pmi}^{-1}(s)f_{pi}(s)\right)
\end{equation}
其中$Z_{pmi}$是第$i$个受位置控制的机械臂阻抗。因此，将式(\ref{eq:9.4.16})与式(\ref{eq:9.4.15})相等得到第$i$个位置控制器
\begin{equation}\label{eq:9.4.17}
a_{pi} = L^{-1}\!\left\{s\left(\dot{x}_{pdi}(s) - Z_{pmi}^{-1}(s)f_{pi}(s)\right)\right\}
\end{equation}
其中$L^{-1}$用于表示逆拉普拉斯变换操作。

继续分离位置控制和力控制设计，我们使用式(\ref{eq:9.4.13})来定义在任务空间方向上受力控制的方程为
\begin{equation}\label{eq:9.4.18}
\ddot{x}_{fj} = a_{fj}
\end{equation}
其中下标$j$表示第$j$个受力控制的任务空间变量，下标$f$表示力控制。与受力控制的任务空间方向相关的环境力用$f_{fj}$表示。

假设零初始条件，式(\ref{eq:9.4.18})的拉普拉斯变换可以写成：
\begin{equation}\label{eq:9.4.19}
s\dot{x}_{fj} = a_{fj}
\end{equation}
从式(\ref{eq:9.4.8})给出的力控制模型，式(\ref{eq:9.4.19})的左侧也可以写成
\begin{equation}\label{eq:9.4.20}
a_{fj} = L^{-1}\!\left\{ s Z_{fmj}^{-1}(s)\bigl(f_{fdj}(s) - f_{fj}(s)\bigr) \right\}
\end{equation}
\begin{equation}\label{eq:9.4.21}
s\dot{x}_{fj}(s) = s Z_{fmj}^{-1}(s)\bigl(f_{fdj}(s) - f_{fj}(s)\bigr)
\end{equation}
其中$Z_{fmj}$是第$j$个受力控制的机械臂阻抗。因此，将式(\ref{eq:9.4.21})与式(\ref{eq:9.4.19})相等得到第$j$个力控制器。

整体的``混合''阻抗控制策略通过将式(\ref{eq:9.4.12})与式(\ref{eq:9.4.17})和式(\ref{eq:9.4.20})一起使用来获得。这种混合阻抗控制策略在表9.4.1中总结。

\textbf{表9.4.1}\quad 混合阻抗控制器

\begin{center}
\small
\begin{tabular}{|l|p{11.2cm}|}
\hline
\textbf{扭矩控制器} & $\displaystyle \tau = M(q)J^{-1}(q)\bigl(a - \dot{J}(q)\dot{q}\bigr) + V_m(q,\dot{q})\dot{q} + G(q) + F(\dot{q}) + J^{\mathsf T}(q)f$，其中$J(q)$为任务空间雅可比矩阵。 \\
\hline
\textbf{位置控制} & $\displaystyle a_{pi} = L^{-1}\!\left\{ s\bigl(\dot{x}_{pdi}(s) - Z_{pmi}^{-1}(s)f_{pi}(s)\bigr) \right\}$ \\
\hline
\textbf{力控制} & $\displaystyle a_{fj} = L^{-1}\!\left\{ s Z_{fmj}^{-1}(s)\bigl(f_{fdj}(s) - f_{fj}(s)\bigr) \right\}$ \\
\hline
\textbf{稳定性} & 对力或位置阶跃输入，稳态误差为零。 \\
\hline
\textbf{注释} & 机械臂阻抗$Z_{pmi}$与$Z_{fmj}$通过对偶原理选取。 \\
\hline
\end{tabular}
\end{center}

注意，高级控制器将用于选择要受位置或力控制的任务空间分量，然后分配适当的机械臂阻抗。如对偶原理所阐述的，式(\ref{eq:9.4.20})中的机械臂阻抗$Z_{fmj}$被分配为非电容性的；用于位置控制的阻抗$Z_{pmi}$被分配为非惯性的。

\textbf{示例 9.4--1：倾斜表面上的混合阻抗控制}

我们希望为图9.2.5给出的机器人机械臂系统构造混合阻抗控制器。忽略关节摩擦与表面摩擦，并假设切向力（即$f_2$）满足关系
\begin{equation}\label{eq:9.4.22}
f_2 = d_e \dot{v},
\end{equation}
法向力（即$f_1$）满足关系
\begin{equation}\label{eq:9.4.23}
f_1 = h_e \ddot{u} + b_e \dot{u} + k_e u,
\end{equation}
其中$h_e$、$b_e$、$d_e$与$k_e$均为正的标量常数。

由表9.4.1，混合阻抗控制器可写为
\begin{equation}\label{eq:9.4.24}
\tau = M(q)J^{-1}(q)\mathbf{a} + G(q) + J^{\mathsf T}(q)f,
\end{equation}
其中$\mathbf{a}$为$2\times 1$向量，表示分离的位置控制与力控制策略；$\tau$、$J$、$G$与$f$的定义与\textbf{示例 9.2--1}中相同。式(\ref{eq:9.4.24})给出的扭矩控制器使机器人动力学在任务空间中解耦为
\begin{equation}\label{eq:9.4.25}
\ddot{\mathbf{z}} = \begin{bmatrix} \ddot{u} \\ \ddot{v} \end{bmatrix} = \mathbf{a},
\end{equation}
其中$\mathbf{z}=\begin{bmatrix} u \\ v \end{bmatrix}$为任务空间坐标向量；因而易于判定哪些任务空间方向应采力控制、哪些应采位置控制。

将定义\ref{def:9.4.2}用于式(\ref{eq:9.4.22})可知，$v$方向上的环境阻抗为电阻性阻抗；因此由对偶原理，选取采用电容性机械臂阻抗的位置控制器：
\begin{equation}\label{eq:9.4.26}
Z_{pm1}(s) = h_m s + b_m + \frac{k_m}{s},
\end{equation}
其中$h_m$、$b_m$与$k_m$均为正的标量常数。由于任务空间变量$v$受位置控制，沿用式(\ref{eq:9.4.14})的记号可得
\[
\ddot{z}_{p1} = \ddot{v} = a_{p1}, \qquad f_{p1} = f_2.
\]
结合式(\ref{eq:9.4.26})与表9.4.1中$a_{p1}$的定义，可直接化为
\begin{equation}\label{eq:9.4.27}
a_{p1} = \ddot{z}_{pd1} + \frac{b_m}{h_m}\bigl(\dot{z}_{pd1} - \dot{z}_{p1}\bigr) + \frac{k_m}{h_m}(z_{pd1} - z_{p1}) - \frac{1}{h_m} f_{p1}.
\end{equation}

将定义\ref{def:9.4.3}用于式(\ref{eq:9.4.23})可知，$u$方向上的环境阻抗为电容性阻抗；因此由对偶原理，选取采用惯性机械臂阻抗的力控制器：
\begin{equation}\label{eq:9.4.28}
Z_{fm1}(s) = d_m s,
\end{equation}
其中$d_m$为正的标量常数。由于任务空间变量$u$受力控制，沿用式(\ref{eq:9.4.18})的记号可得
\[
\ddot{z}_{f1} = \ddot{u} = a_{f1}, \qquad f_{f1} = f_1.
\]
结合式(\ref{eq:9.4.28})与表9.4.1中$a_{fj}$的定义，可直接化为
\begin{equation}\label{eq:9.4.29}
a_{f1} = \frac{1}{d_m}(f_{fd1} - f_{f1}).
\end{equation}

整体的阻抗控制策略可将
\begin{equation}\label{eq:9.4.30}
\mathbf{a} = \begin{bmatrix} a_{f1} \\ a_{p1} \end{bmatrix}
\end{equation}
代入式(\ref{eq:9.4.24})得到，其中$a_{p1}$与$a_{f1}$分别由式(\ref{eq:9.4.27})与式(\ref{eq:9.4.29})给出。

%-------------------------------------------
\subsection{实现要点}
%-------------------------------------------

如前所述，机械臂阻抗$Z_m(s)$的选择应使对偶原理得以满足。然而从实现角度看，还应在$a_{pi}$与$a_{fj}$的表达式中选取$Z_m(s)$，使得控制器仅需$f$、$x$与$\dot{x}$的测量，而不依赖加速度$\ddot{x}$或力的导数$\dot{f}$（二者在实际中往往难以可靠得到）。若将机械臂阻抗取为
\begin{equation}\label{eq:9.4.31}
Z_m(s) = hs + Z_r(s),
\end{equation}
其中$h$为某一正的标量常数，且$Z_r(s)$选为真有理传递函数[Anderson and Spong 1988]，则可以证明：在上述阻抗选取下，$a_{pi}$与$a_{fj}$的表达式中将不再需要$\ddot{x}$与$\dot{f}$的测量。

正如表9.4.1所概括，混合阻抗控制器需要测量任务空间力。与混合位置/力控制器情形类似，可通过坐标变换由传感器力得到任意应用场景下的任务空间力。具体而言，传感器力与任务空间力满足
\begin{equation}\label{eq:9.4.32}
f = J^{-\mathsf T}(q) J_s^{\mathsf T}(q) f_s,
\end{equation}
其中$J_s(q)$为$n\times n$传感器雅可比矩阵，$f_s$为$n\times 1$传感器力向量。

在上述实现考量之下，整体混合阻抗控制策略如图9.4.4所示。框图中的前馈项用于表示控制律中的项
\begin{equation}\label{eq:9.4.33}
-M(q)J^{-1}(q)\dot{J}(q)\dot{q} + V_m(q,\dot{q})\dot{q} + G(q) + F(\dot{q}) + J^{\mathsf T}(q)f.
\end{equation}

%===========================================
\section{降维状态位置/力控制}\label{sec:9.5}
%===========================================

如果一个刚性机械臂受到刚性环境的约束，自由度就会降低，因为机械臂末端执行器不能穿过环境移动；因此，关于位置的一个或多个自由度会丢失。因此，当机械臂末端执行器接触环境约束时，末端执行器和环境之间会产生相互作用力（有时称为约束力）。这种``减少''位置自由度同时发展约束力的过程使人相信应该根据这一自然现象设计位置/力控制器。

最近，研究人员开始制定一个理论框架[McClamroch and Wang 1988]、[Kankaanranta and Koivo 1988]，利用动力学中的经典结果将约束力的影响纳入机器人机械臂模型。假设应该为接触刚性约束的刚性机械臂设计控制器的理由是，在大多数力控制应用中，环境比机械臂刚硬得多。因此，假设机械臂是刚性的而环境是柔性的似乎不合理。

%-------------------------------------------
\subsection{完整约束对机械臂动力学的影响}
%-------------------------------------------

对于本节给出的控制器，我们假设环境约束是完整且无摩擦的。也就是说，我们假设存在一个关节空间坐标中的约束函数$\bar{\psi}(q)$（一个$p \times 1$向量函数）满足
\begin{equation}\label{eq:9.5.1}
\bar{\psi}(q) = 0
\end{equation}

式(\ref{eq:9.5.1})给出的关系说明环境约束是完整的。假设约束函数的维度小于关节变量的数量（即$p < n$）。对于特定问题，从机器人运动学和环境配置中找到函数$\bar{\psi}(q)$。为了阐明完整概念，我们现在给出一个示例。

\textbf{示例 9.5--1：完整约束}

我们希望为图9.2--5和9.2--6中给出的机器人/环境配置制定约束函数$\bar{\psi}(q)$。在这两种配置中，关节空间的维度为2，约束函数的维度为1（即，约束是一维表面）。对于图9.2.5给出的机械臂系统，约束函数可取为
\begin{equation*}
\bar{\psi}(q) = q_1 - q_2 - 3 = 0.
\end{equation*}

对于图9.2.6给出的机械臂系统，约束函数由下式给出
\begin{equation*}
\bar{\psi}(q) = \frac{1}{4}q_1^2 + q_2^2 - 1 = 0.
\end{equation*}

对于具有完整且无摩擦约束的通用$n$连杆机器人机械臂，约束机器人动力学可以写成形式
\begin{equation}\label{eq:9.5.2}
\tau = M(q)\ddot{q} + V_m(q, \dot{q})\dot{q} + G(q) + F(\dot{q}) + A^{\mathsf T}(q)\lambda
\end{equation}
其中$\lambda$是一个$p \times 1$向量，表示与约束相关的广义力乘数（拉格朗日乘子），约束雅可比矩阵$A(q)$是一个$p \times n$矩阵，定义为
\begin{equation}\label{eq:9.5.3}
A(q) = \frac{\partial \bar{\psi}(q)}{\partial q}
\end{equation}

如[Kankaanranta and Koivo 1988]中所述，我们假设$A^{\mathsf T}(q)$的$p$列在关节空间上线性独立。此外，假定拉格朗日乘子向量$\lambda$不依赖于$q$与$\dot{q}$。

为了激励式(\ref{eq:9.5.2})的来源，我们重新检查拉格朗日方程
\begin{equation}\label{eq:9.5.4}
\frac{d}{dt}\left[\frac{\partial L}{\partial \dot{q}}\right] - \frac{\partial L}{\partial q} = \tau
\end{equation}
其中受约束机器人机械臂的修正拉格朗日量由下式给出
\begin{equation}\label{eq:9.5.5}
L = K - P - \lambda^{\mathsf T} \bar{\psi}(q)
\end{equation}
注意，式(\ref{eq:9.5.5})中给出的拉格朗日量实际上与第2章中给出的拉格朗日量相同，因为
\[
\bar{\psi}(q) = 0
\]
如式(\ref{eq:9.5.1})所要求的。

从式(\ref{eq:9.5.4})和式(\ref{eq:9.5.5})中，我们可以看到机器人机械臂动力学方程的结构对于受约束的机械臂与对于不受约束的机械臂相同，除了拉格朗日方程中替换$\bar{\psi}(q) = 0$所产生的任何新项之外。具体来说，式(\ref{eq:9.5.2})给出的约束机器人机械臂方程是通过将式(\ref{eq:9.5.5})代入式(\ref{eq:9.5.4})并利用恒等式
\[
\frac{\partial}{\partial q}\bigl(\lambda^{\mathsf T}\bar{\psi}(q)\bigr) = \biggl[\frac{\partial \bar{\psi}(q)}{\partial q}\biggr]^{\mathsf T}\lambda = A^{\mathsf T}(q)\lambda
\]
得到的。

在我们开发降维状态位置/力控制器之前，我们简要地从启发式的角度解释位置和力变量如何被降维。首先，注意对于式(\ref{eq:9.5.2})给出的约束机器人动力学，变量$\lambda$用于表示应该被控制的约束力。注意，$\lambda$的维度是$p$，根据定义，小于$n$。在前面给出的环境力公式中，力的维度被假定为等于$n$。在这种方法中，我们能够减少必须控制的力的数量，因为我们最初假设约束表面是无摩擦的。当然，在现实中，表面摩擦将存在于大多数机器人力控制应用中；然而，它通常被视为干扰，因为这样的摩擦是所施加的法向接触力的函数。

其次，我们检查式(\ref{eq:9.5.1})给出的机器人运动的位置约束。因为根据假设$\bar{\psi}(q) = 0$，我们可以对式(\ref{eq:9.5.1})关于时间求导得到
\begin{equation}\label{eq:9.5.6}
A(q)\dot{q} = 0
\end{equation}
其中$A(q)$在式(\ref{eq:9.5.3})中定义。式(\ref{eq:9.5.6})给出的表达式给出了运动学速度约束的简洁形式。从式(\ref{eq:9.5.1})和式(\ref{eq:9.5.6})给出的位置和速度约束，我们可以说明机械臂动力学属于在$\mathbb{R}^{2n}$中由$C$定义的不变流形
\[
C = \{(q, \dot{q}) : \bar{\psi}(q) = 0,\ A(q)\dot{q} = 0\}
\]
也就是说，机器人机械臂的运动保持在由$C$定义的流形上。如[McClamroch and Wang 1988]中所述，流形$C$在$\mathbb{R}^{2n}$上是奇异的；因此，我们可以降低运动动力学的阶数。

%-------------------------------------------
\subsection{降维建模和控制}
%-------------------------------------------

由单个$n$关节非冗余机械臂受刚性环境约束组成的机器人系统具有$n - p$个运动自由度。注意，然而，如式(\ref{eq:9.5.2})给出的机械臂关节变量模型包含$n$个位置变量$(q)$，结合$p$个力变量$(\lambda)$，使控制变量（即状态）的总数$n + p$超过控制输入的数量$n$。在本节中，我们展示了如何使用变量变换来降低动力学模型的状态，从而将控制变量的数量从$n + p$减少到$n$。

降维模型可以通过用另一组独立坐标（称为约束空间坐标）表示式(\ref{eq:9.5.2})给出的机械臂动力学来获得，这组坐标将用$x$表示。预期$x$是一个$(n - p) \times 1$关节空间坐标向量。也就是说，$x$是$q$的一个子集。

通过假设存在一个将约束空间向量$x$[一个$(n - p) \times 1$向量]与关节空间向量$q$相关联的$n \times 1$向量函数$g(x)$来降低约束机器人模型。该函数由下式给出
\begin{equation}\label{eq:9.5.7}
q = g(x)
\end{equation}
其中函数$g(x)$必须被选择为使得
\begin{equation}\label{eq:9.5.8}
\left\lvert\left[\frac{\partial g(x)}{\partial x}\right]^{\mathsf T} A^{\mathsf T}(q)\right\rvert_{\,q=g(x)} = 0,
\end{equation}
其中竖线记号与``矩阵$\bigl[\frac{\partial g(x)}{\partial x}\bigr]^{\mathsf T} A^{\mathsf T}(q)$在$q=g(x)$处恒为零''等价（与式(\ref{eq:9.5.9})中$\Sigma^{\mathsf T}$写法一致）；当$n-p=p$且括号内为方阵时，上式即该方阵的行列式为零。
并且$n \times (n - p)$雅可比矩阵$\Sigma(x)$定义为
\begin{equation}\label{eq:9.5.9}
\Sigma(x) = \frac{\partial g(x)}{\partial x}
\end{equation}
在由$x$给出的约束空间运动中包含$n - p$个独立行。

即使约束空间向量$x$的维度小于关节空间向量$q$的维度，通常也能够找到式(\ref{eq:9.5.7})中给出的从$x$到$q$的函数映射。对于特定问题，式(\ref{eq:9.5.1})给出的完整方程和机器人运动学的代数关系用于找到$g(x)$。还要注意，$g(x)$的选择是非唯一的，并且式(\ref{eq:9.5.8})中给出的$g(x)$条件以及$\Sigma(x)$行的条件是相互关联的。具体来说，$\Sigma(x)$包含$n - p$个独立行的条件确保了解耦模型表示$n - p$个独立方程。式(\ref{eq:9.5.8})中给出的$g(x)$条件用于确保由$\lambda$表示的力可以与由$x$表示的约束空间运动解耦。应该注意的是，由于约束被假定为完整的，并且矩阵$A^T(q)$被假定为具有$p$个线性独立的列，我们可以证明上述两个条件始终成立。读者可以参考[McClamroch and Wang 1988]了解详情。

为了获得用约束空间坐标表示的降维动力学，我们首先对式(\ref{eq:9.5.7})关于时间求导得到
\begin{equation}\label{eq:9.5.10}
\dot{q} = \Sigma(x)\dot{x}
\end{equation}
其中$\Sigma(x)$在式(\ref{eq:9.5.9})中定义。对式(\ref{eq:9.5.10})关于时间求导得到
\begin{equation}\label{eq:9.5.11}
\ddot{q} = \Sigma(x)\ddot{x} + \dot{\Sigma}(x)\dot{x}
\end{equation}

在将式(\ref{eq:9.5.7})、式(\ref{eq:9.5.10})和式(\ref{eq:9.5.11})中的$q$、$\dot{q}$和$\ddot{q}$分别代入式(\ref{eq:9.5.2})后，我们得到降维模型
\begin{equation}\label{eq:9.5.12}
M(x)\Sigma(x)\ddot{x} + N(x, \dot{x}) + A^T(x)\lambda = \tau
\end{equation}
其中
\[
N(x, \dot{x}) = (V_m(x, \dot{x})\Sigma(x) + M(x)\dot{\Sigma}(x))\dot{x} + G(x) + F(x)
\]

式(\ref{eq:9.5.12})给出的降维模型的意义通过用$\Sigma^{\mathsf T}(x)$前乘式(\ref{eq:9.5.12})来说明
\begin{equation}\label{eq:9.5.13}
\tau^* = M^*\ddot{x} + N^*,
\end{equation}
其中
\[
\tau^* = \Sigma^{\mathsf T}(x)\tau, \quad M^* = \Sigma^{\mathsf T}(x)M(x)\Sigma(x), \quad N^* = \Sigma^{\mathsf T}(x)N(x, \dot{x}).
\]
注意，式(\ref{eq:9.5.12})中的接触力已被式(\ref{eq:9.5.8})的结果移除，式(\ref{eq:9.5.8})与$A(q)\Sigma(x)=0$（在$q=g(x)$上取值）等价。式(\ref{eq:9.5.13})给出的模型是有用的，因为控制机械臂在约束表面上运动的动力学已经从$n$个微分方程减少到$n - p$个微分方程；此外，运动已经与接触力解耦。这两个结果在后续位置/力控制器的设计和稳定性分析中很重要。

在介绍降维状态位置/力控制器之前，我们给出一些关于位置/力跟踪问题的定义。首先，约束位置跟踪误差定义为
\begin{equation}\label{eq:9.5.14}
\tilde{x} = x_d - x
\end{equation}
我们假设期望的约束空间轨迹及其前两个导数，分别用$x_d$、$\dot{x}_d$和$\ddot{x}_d$表示，都是有界函数。我们还假设期望的力乘数轨迹$\lambda_d$是一个已知的界函数，由此定义了相应的力乘数跟踪误差变量
\begin{equation}\label{eq:9.5.15}
\tilde{\lambda} = \lambda_d - \lambda
\end{equation}

降维状态位置/力控制器是一个反馈线性化控制器；因此，需要精确了解机器人动力学。降维状态控制器[McClamroch and Wang 1988]是
\begin{equation}\label{eq:9.5.16}
\tau = M(x)\Sigma(x)(\ddot{x}_d + K_v\dot{\tilde{x}} + K_p \tilde{x}) + N(x, \dot{x}) + A^{\mathsf T}(x)(\lambda_d + K_f \tilde{\lambda})
\end{equation}
其中$K_v$和$K_p$是对角、正定的$(n - p) \times (n - p)$矩阵，$K_f$是对角、正定的$p \times p$矩阵。

为了确定位置误差和力跟踪误差的稳定性类型，我们将式(\ref{eq:9.5.16})代入式(\ref{eq:9.5.12})得到
\begin{equation}\label{eq:9.5.17}
M(x)\Sigma(x)(\ddot{\tilde{x}} + K_v\dot{\tilde{x}} + K_p \tilde{x}) + A^{\mathsf T}(x)(\tilde{\lambda} + K_f \tilde{\lambda}) = 0
\end{equation}

用$\Sigma^{\mathsf T}(x)$前乘式(\ref{eq:9.5.17})得到
\begin{equation}\label{eq:9.5.18}
M^*(\ddot{\tilde{x}} + K_v\dot{\tilde{x}} + K_p \tilde{x}) = 0
\end{equation}
作为式(\ref{eq:9.5.8})的结果。因为$\Sigma^{\mathsf T}(x)$在约束空间运动中包含$n - p$个独立列，$M(x)$是一个正定对称矩阵，$M^*$是一个正定对称矩阵。因此，我们可以用$M^{*-1}$前乘式(\ref{eq:9.5.18})得到
\begin{equation}\label{eq:9.5.19}
\ddot{\tilde{x}} + K_v\dot{\tilde{x}} + K_p \tilde{x} = 0
\end{equation}

因为$K_v$和$K_p$是对角正定矩阵，可以对式(\ref{eq:9.5.19})应用标准线性控制论证，得到
\begin{equation}\label{eq:9.5.20}
\lim_{t \to \infty} \ddot{\tilde{x}} = 0, \quad \lim_{t \to \infty} \dot{\tilde{x}} = 0, \quad \lim_{t \to \infty} \tilde{x} = 0
\end{equation}

为了获得力跟踪误差的稳定性结果，我们将式(\ref{eq:9.5.20})代入式(\ref{eq:9.5.17})得到
\begin{equation}\label{eq:9.5.21}
\lim_{t \to \infty} A^{\mathsf T}(x)(I + K_f)\tilde{\lambda} = 0
\end{equation}
其中式(\ref{eq:9.5.21})中的$I$是$p \times p$单位矩阵。因为$A^{\mathsf T}(x)$的$p$列被假定为线性独立的，复合矩阵$(I + K_f)$是一个正定对角矩阵，我们可以将式(\ref{eq:9.5.21})写成
\begin{equation}\label{eq:9.5.22}
\lim_{t \to \infty} \tilde{\lambda} = 0
\end{equation}

如式(\ref{eq:9.5.20})和式(\ref{eq:9.5.22})所阐述的，式(\ref{eq:9.5.16})给出的降维状态位置/力控制器对约束位置误差和力跟踪误差都产生了渐近稳定性结果。降维状态位置/力控制策略在表9.5.1中总结。

\textbf{表9.5.1}\quad 降维状态位置/力控制器

\begin{center}
\small
\begin{tabular}{|l|p{11.2cm}|}
\hline
\textbf{扭矩控制器} & \begin{minipage}[t]{\linewidth}\raggedright\setlength{\lineskip}{6pt}
$\displaystyle \tau = M(x)\Sigma(x)(\ddot{x}_d + K_v\dot{\tilde{x}} + K_p \tilde{x}) + N(x,\dot{x}) + A^{\mathsf T}(x)(\lambda_d + K_f \tilde{\lambda})$，其中$A^{\mathsf T}(x)$即约束雅可比的转置；\\
$\displaystyle N(x,\dot{x}) = \bigl(V(x,\dot{x})\Sigma(x) + M(x)\dot{\Sigma}(x)\bigr)\dot{x} + G(x) + F(x,\dot{x})$（此处$V$与离心/Coriolis矩阵记号对应正文中的$V_m$）。
\par\end{minipage} \\
\hline
\textbf{稳定性} &
\setlength{\parskip}{6pt}
\textbf{(i)}\ 位置跟踪（约束坐标）：$\displaystyle \lim_{t\to\infty} x(t) = x_d(t)$。\par
\textbf{(ii)}\ 力跟踪（乘子）：$\displaystyle \lim_{t\to\infty} \lambda(t) = \lambda_d(t)$。 \\
\hline
\textbf{注释} &
\setlength{\parskip}{6pt}
假定约束为完整且无摩擦。\par
为使位置/力解耦，$A^{\mathsf T}(q)$在关节空间中必须包含$p$个线性独立的列。 \\
\hline
\end{tabular}
\end{center}

如表9.5.1所述，降维状态位置/力控制器需要测量力变量$\lambda$。与混合位置/力情形类似，可由传感器力经坐标变换得到$\lambda$。具体而言，$\lambda$与传感器力满足
\begin{equation}\label{eq:9.5.23}
A^{\mathsf T}(q)\lambda = J_s^{\mathsf T}(q) f_s,
\end{equation}
其中$J_s(q)$为$n\times n$传感器雅可比矩阵，$f_s$为$n\times 1$传感器力向量。由式(\ref{eq:9.5.23})可得到$p$个独立方程，从而将$\lambda$表示为仅依赖于可测$q$与$f_s$的函数。

\textbf{示例 9.5--2：倾斜表面上的降维状态位置/力控制}

我们希望为图9.2.5给出的机器人机械臂系统构造降维状态位置/力控制器。忽略关节摩擦与表面摩擦。由示例~9.5--1，该问题的完整约束可取为
\[
\bar{\psi}(q) = q_1 - q_2 - 3 = 0.
\]
利用式(\ref{eq:9.5.2})与式(\ref{eq:9.5.3})，约束面上的机器人动力学可写为
\begin{equation}\label{eq:9.5.24}
\tau = M\ddot{q} + G + A^{\mathsf T}\lambda,
\end{equation}
其中$\tau$、$M$、$q$与$G$的定义与\textbf{示例 9.2--1}相同；又因$p=1$，可将$A^{\mathsf T}$取为
\begin{equation}\label{eq:9.5.25}
A^{\mathsf T} = \begin{bmatrix} 1 \\ -1 \end{bmatrix}.
\end{equation}
选取降维坐标$x=q_1$，并要求存在$q=g(x)$使式(\ref{eq:9.5.7})成立；由上述约束直接得到
\begin{equation}\label{eq:9.5.26}
\begin{bmatrix} q_1 \\ q_2 \end{bmatrix} = \begin{bmatrix} x \\ x - 3 \end{bmatrix}.
\end{equation}
利用式(\ref{eq:9.5.9})，
\begin{equation}\label{eq:9.5.27}
\Sigma = \begin{bmatrix} 1 \\ 1 \end{bmatrix}.
\end{equation}
由表9.5.1，扭矩控制器为
\begin{equation}\label{eq:9.5.28}
\tau = M\Sigma(\ddot{x}_d + K_v\dot{\tilde{x}} + K_p \tilde{x}) + G + A^{\mathsf T}(\lambda_d + K_f \tilde{\lambda}),
\end{equation}
其中$x_d$、$\dot{x}_d$与$\ddot{x}_d$分别为$x=q_1$的期望轨迹及其导数，$\tilde{x}=x_d-x$，$\dot{\tilde{x}}=\dot{x}_d-\dot{x}$；$\lambda_d$为期望力乘子，$\tilde{\lambda}=\lambda_d-\lambda$。

定义$x_d$与$\lambda_d$的一种做法是先把期望的任务空间法向力$f_{d1}$与沿表面的期望位置$v_d$联系起来。由\textbf{示例 9.2--1}中的雅可比与力的记号，
\begin{equation}\label{eq:9.5.29}
A^{\mathsf T}\lambda = J^{\mathsf T}f.
\end{equation}
忽略表面摩擦（即$f_2=0$），由上式可得标量关系
\begin{equation}\label{eq:9.5.30}
\lambda = \frac{f_1}{\sqrt{2}}.
\end{equation}
此外，由$h(q)$的第二分量，切向坐标$v=\frac{1}{\sqrt{2}}(q_1+q_2)$；并结合$\bar{\psi}(q)=0$，可写成
\begin{equation}\label{eq:9.5.31}
q_1 = \frac{\sqrt{2}}{2}\,v + \frac{3}{2}.
\end{equation}
因此，可把$f_{d1}$代入式(\ref{eq:9.5.30})得到$\lambda_d$，并把$v_d$代入式(\ref{eq:9.5.31})得到相应的$q_{d1}$，进而用它生成$x_d=q_{d1}$及其导数。

%-------------------------------------------
\subsection{实现要点}
%-------------------------------------------

考虑到上述实现结构，整体降维位置/力控制策略可用与此类似的框图表示（参见原著插图）。框图中的前馈项对应于
\begin{equation}\label{eq:9.5.32}
\bigl(V_m(x,\dot{x})\Sigma(x) + M(x)\dot{\Sigma}(x)\bigr)\dot{x} + G(x) + F(x,\dot{x}).
\end{equation}

%===========================================
\section{总结}\label{sec:9.6}
%===========================================

在本章中，给出了几种用于刚性机器人的位置/力控制策略。目的是按时间顺序记录力控制的发展。一些新的研究领域，如执行器动力学、传感器动力学、关节柔性、机械臂动态不确定性、表面不确定性和环境影响对位置/力控制器性能的影响，现在正在被许多研究人员研究。在解决了其中一些问题之后，可以期望看到机器人机械臂更频繁地用于力控制应用。

%===========================================
% 参考文献
%===========================================
\begin{thebibliography}{99}
\bibitem{anderson1988} Anderson, R., and M. Spong, ``Hybrid impedance control of robotic manipulators,'' \textit{J. Robot. Autom.}, vol. 4, no. 5, pp. 549--556, Oct. 1988.

\bibitem{chae1988} Chae, A., C. Atkeson, and J. Hollerbach, \textit{Model-Based Control of a Robot Manipulator}. Cambridge, MA: MIT Press, 1988.

\bibitem{hogan1987} Hogan, N., ``Stable execution of contact tasks using impedance control,'' \textit{Proc. IEEE Int. Conf Robot. Autom.}, pp. 595--601, Raleigh NC, Mar. 1987.

\bibitem{kankaanranta1988} Kankaanranta, R., and H. Koivo, ``Dynamics and simulation of compliant motion of a manipulator,'' \textit{IEEE Trans. Robot. Autom.}, vol. 4, pp. 163--173, Apr. 1988.

\bibitem{lipkin1988} Lipkin, H., and J. Duffy, ``Hybrid twist and wrench control for a robot manipulator,'' \textit{Trans. ASME J. Mech. Transmissions Autom. Design}, vol. 110, pp. 138--144, June 1988.

\bibitem{mcclamroch1988} McClamroch, N., and D. Wang, ``Feedback stabilization and tracking of constrained robots,'' \textit{IEEE Trans. Autom. Control}, vol. 33, no. 5, pp. 419--426, May 1988.

\bibitem{raibert1981} Raibert, M., and J. Craig, ``Hybrid position/force control of manipulators,'' \textit{J. Dyn. Syst. Meas. Control}, vol. 102, pp. 126--132, June 1981.

\bibitem{salisbury1980} Salisbury, J., and J. Craig, ``Active stiffness control of manipulator in Cartesian coordinates,'' \textit{Proc. 19th IEEE Conf. Decision Control}, Dec. 1980.
\end{thebibliography}

%===========================================
% 习题
%===========================================
\section*{习题}

\subsection*{第9.2节}

\textbf{9.2--1} 对于图9.2.6中给出的机械臂，使用新的表面函数找到施加在环境表面上的法向力和切向力相关的任务空间雅可比矩阵
\[
\frac{1}{4}q_1^2 + q_2^2 = 1
\]

\textbf{9.2--2} 为图9.2.5中给出的机器人机械臂系统设计并模拟一个刚度控制器，表面函数由下式给出
\[
\frac{1}{2}q_1 - q_2 = 5
\]
控制目标是将末端执行器移动到$q_{d2} = 0.3$ m的期望最终位置，同时施加$2$ N的最终期望法向力。忽略表面和关节摩擦，并假设法向力满足关系
\[
f_1 = k_e(u - u_e)
\]
其中$k_e = 100$ N/m，$u_e$是到表面的静态法向距离。机器人连杆质量可以假定为单位质量，而初始末端执行器位置被假定为由下式给出
\[
v = \frac{1}{\sqrt{2}}(q_1 + q_2) = 0
\]

\subsection*{第9.3节}

\textbf{9.3--1} 由于混合位置/力控制策略的位置和力误差系统是线性的，讨论如何选择控制器增益来改变机械臂末端执行器的性能。

\textbf{9.3--2} 为图9.2.5中给出的机器人机械臂系统设计并模拟一个混合位置/力控制器，表面函数由下式给出
\[
\frac{1}{\sqrt{2}}q_1 + \frac{1}{\sqrt{2}}q_2 - 5 = 0
\]
控制目标是使末端执行器沿表面轨迹移动
\[
v_d = \sin(t) \text{ m}
\]
同时施加期望法向力
\[
f_{d1} = 1 + e^{-t} \text{ N}
\]
忽略表面和关节摩擦，并假设法向力满足关系
\[
f_1 = k_e(u - u_e)
\]
其中$k_e = 1000$ N/m，$u_e$是到表面的静态法向距离。机器人连杆质量可以假定为单位质量，而初始末端执行器位置被假定为由下式给出
\[
u = \frac{1}{\sqrt{2}}(q_1 - q_2) = 0
\]

\textbf{9.3--3} 解释为什么混合位置/力控制策略实际上是一种位置控制策略。

\subsection*{第9.4节}

\textbf{9.4--1} 假设对于力或位置控制，机械臂阻抗被选择为
\[
Z_m(s) = hs + Z_r(s)
\]
其中$h$为正常数，$Z_r(s)$是一个适当稳定的传递函数。使用上述机械臂阻抗，证明混合阻抗控制器只需要测量关节位置、关节速度和接触力。

\textbf{9.4--2} 为图9.2.5中给出的机器人机械臂系统设计一个混合阻抗控制器，表面函数由下式给出
\[
\frac{1}{\sqrt{2}}q_1 + \frac{1}{\sqrt{2}}q_2 - \frac{3}{\sqrt{2}} = 0
\]
忽略表面和关节摩擦，并假设法向力和切向力满足关系
\[
f_1 = k_e(u - u_e) + b_e \dot{u}
\]
其中$k_e = 10$ N/m，$b_e = 1$ N-s/m，$u_e$是到表面的静态法向距离。

\textbf{9.4--3} 明确展示如何使用表9.4.1得到示例~9.4--1中式(\ref{eq:9.4.27})给出的$a_{p1}$表达式。

\subsection*{第9.5节}

\textbf{9.5--1} 为图9.2.6中给出的机器人机械臂系统设计一个降维状态控制器。

\textbf{9.5--2} 对于在问题9.5--1中开发的控制器，展示法向力（即$f_1$）如何与力乘数（即$\lambda$）相关，以及降维运动变量（即$x$）如何与沿表面的运动（即$v$）相关。

\ifstandalone
  \end{document}
\fi
